\section{Bilateral conservation and arithmetic--geometric transduction}
\label{lib04:bilateral}

The holographic constant of arithmetic--geometric coupling denotes
the generated record \(K\), with its ordered windows, phase origin,
orientation, and incidences. Its production through the signed record and
Hadamard inversion is developed in \cref{sec:k-registro}.
The rational reading \(\kappa\), the orbital frame \(O_K\), and the
operators transporting its realizations are different objects.
This distinction allows their functions to be composed: generation of the record,
arithmetic encoding, information transport, and recovery of
geometric readings.

The regional coordinates and \(\alpha\) are supplied by the
APP--TRIT--TPK genealogy, by the same enriched state and its joint
discrete structure of the continuum. The following constructions act on those
subsequent realizations. Their additional content is dimension-independent
bilateral conservation, a reversible extension admitting
incoming memory, and an archiving policy with uniformly
controlled recovery.

\subsection{Bilateral law for two isometric transports}
\label{lib04:ley}

Let \(H,H'\) be Hilbert spaces and
\(\mathcal B,\mathcal C:H\to H'\) linear isometries, so that
\(\mathcal B^*\mathcal B=\mathcal C^*\mathcal C=I_H\).
Calligraphic type distinguishes these transports from the arithmetic
base of refinement. For \(a>0\), define
\begin{equation}
 \begin{aligned}
 p&=a^2+a+1=a(a+1)+1,\\
 N&=(2a+1)p=(a+1)^3+a^3,\\
 Q_-&=a\mathcal B-\mathcal C,\qquad
 Q_+=(a+1)\mathcal B+\mathcal C.
 \end{aligned}
 \label{lib04:def-q}
\end{equation}

\begin{theorem}[Balance of two isometric realizations]
\label{lib04:teo-balance}
For the maps \eqref{lib04:def-q},
\begin{equation}
 (a+1)Q_-^*Q_-+aQ_+^*Q_+=NI_H.
 \label{lib04:balance}
\end{equation}
The identity holds in every dimension, including infinite-dimensional
Hilbert spaces, and admits rectangular transports with a common
domain and codomain.
\end{theorem}

\begin{proof}
Write \(X=\mathcal B^*\mathcal C+\mathcal C^*\mathcal B\).
The isometric property of both transports gives
\begin{align*}
 Q_-^*Q_-&=(a^2+1)I-aX,\\
 Q_+^*Q_+&=((a+1)^2+1)I+(a+1)X.
\end{align*}
The weights \(a+1\) and \(a\) cancel the terms involving
\(X\) exactly. The remaining coefficient is
\[
 (a+1)(a^2+1)+a((a+1)^2+1)
 =2a^3+3a^2+3a+1=N.
\]
The compositions involve bounded operators, and the calculation uses
only their products and adjoints. The proof therefore also applies
in infinite dimension.
\end{proof}

For \(a=8\), supplied by the preceding nonadic partition,
\begin{equation}
 9\|Q_-v\|^2+8\|Q_+v\|^2
 =(9^3+8^3)\|v\|^2=17\cdot73\,\|v\|^2.
 \label{lib04:balance-canonico}
\end{equation}
The factor \(73=8\cdot9+1\) normalizes this general identity. In the
ternary module it is also realized as a determinant and a similarity
factor; each of these roles retains its own hypotheses.
The balance \eqref{lib04:balance} uses the isometries and admits both
finite-order transports and transports without that property.

\subsection{Factorization, recovery, and channel compatibility}
\label{lib04:factorizacion}

Define
\begin{equation}
 E_av=\frac1{\sqrt p}
 \begin{pmatrix}\sqrt{a(a+1)}\,\mathcal Bv\\\mathcal Cv\end{pmatrix},
 \qquad
 R_a=\frac1{\sqrt{2a+1}}
 \begin{pmatrix}\sqrt a&-\sqrt{a+1}\\
                 \sqrt{a+1}&\sqrt a\end{pmatrix}.
 \label{lib04:factorizacion-er}
\end{equation}
The matrix \(R_a\) acts on the two channels and is understood to be
tensored with \(I_{H'}\). One verifies
\(E_a^*E_a=I_H\), \(R_a^*R_a=I_{H'\oplus H'}\), and
\begin{equation}
 W_av=R_aE_av
 =\frac1{\sqrt N}
 \begin{pmatrix}\sqrt{a+1}\,Q_-v\\\sqrt a\,Q_+v\end{pmatrix},
 \qquad W_a^*W_a=I_H.
 \label{lib04:w}
\end{equation}
When \(a=8\), the initial allocation has weights \(72/73\) and
\(1/73\); the second operation is a rotation between the channels.
This factorization expresses the role of \(72+1\) within the total
norm, distinguishing it from the linear restoration in
\eqref{lib03:cadena-restauracion}.

\begin{proposition}[Bilateral recovery]
\label{lib04:prop-recuperacion}
The unnormalized channels \(x=Q_-v\), \(y=Q_+v\) determine
\begin{equation}
 \mathcal Bv=\frac{x+y}{2a+1},\qquad
 \mathcal Cv=\frac{ay-(a+1)x}{2a+1},\qquad
 v=\mathcal B^*\frac{x+y}{2a+1}.
 \label{lib04:inversa-canales}
\end{equation}
From the normalized channels one recovers
\(v=W_a^*W_av\), with \(\|W_a^*\|=1\) when \(H\neq\{0\}\).
\end{proposition}

\begin{proof}
Adding the definitions of \(Q_\pm\) eliminates \(\mathcal Cv\).
The combination \(aQ_+-(a+1)Q_-\) eliminates \(\mathcal Bv\).
The first isometry has left inverse \(\mathcal B^*\), which
proves the formulas. The final assertion follows from
\eqref{lib04:w}: the adjoint of an isometry has norm one.
\end{proof}

The image \(\operatorname{im}W_a\) is closed, and its orthogonal
projector is \(W_aW_a^*\). A normalized pair is compatible with a
unique input exactly when it belongs to that image; its orthogonal
component quantifies the compatibility defect.

If \(\mathcal B=I\) and \(\mathcal C=U\) is unitary, the two
\(Q_\pm\) are polynomials in the same operator. For the unnormalized
channels, the criterion simplifies to
\begin{equation}
 Q_+x=Q_-y.
 \label{lib04:compatibilidad-unitaria}
\end{equation}
Indeed, this equality is equivalent to
\(U(x+y)=ay-(a+1)x\). Defining
\(v=(x+y)/(2a+1)\), substitution gives
\(Q_-v=x\), \(Q_+v=y\). Necessity follows from commutation
of the two polynomials. This simplified form belongs to the stated
unitary domain; the image criterion covers general transports
between different spaces.

\subsection{Transformation of readers and relative symmetry}
\label{lib04:lectores}

\begin{theorem}[A common reader in both channels]
\label{lib04:teo-lector}
For every bounded self-adjoint operator \(F\) on \(H'\),
\begin{equation}
 W_a^*\operatorname{diag}(F,F)W_a
 =\frac{a(a+1)\mathcal B^*F\mathcal B+
                  \mathcal C^*F\mathcal C}{a(a+1)+1}.
 \label{lib04:lector-general}
\end{equation}
If \(\mathcal B\) is unitary and \(\mathcal C=\mathcal BR\), with
\(R\) unitary, the expression, for \(G=\mathcal B^*F\mathcal B\),
is
\begin{equation}
 \mathcal T_a(G)=\frac{a(a+1)G+R^*GR}{a(a+1)+1}.
 \label{lib04:lector-relativo}
\end{equation}
One has \(\mathcal T_a(G)=G\) if and only if \(GR=RG\).
\end{theorem}

\begin{proof}
Expanding \((a+1)Q_-^*FQ_-+aQ_+^*FQ_+\) cancels the
terms \(\mathcal B^*F\mathcal C+\mathcal C^*F\mathcal B\).
The coefficients of the remaining terms are
\(a(a+1)(2a+1)\) for \(\mathcal B^*F\mathcal B\) and \(2a+1\)
for \(\mathcal C^*F\mathcal C\). Division by \(N\) proves
\eqref{lib04:lector-general}. Substituting \(\mathcal C=\mathcal BR\)
gives \eqref{lib04:lector-relativo}. Equality with \(G\) is equivalent
to \(R^*GR=G\), and unitarity of \(R\) makes this equivalent to
\(GR=RG\).
\end{proof}

The reader \(F=I\) recovers the norm. For a general reader, the
identity determines its transformation, and the coincidence
\(\mathcal B^*F\mathcal B=\mathcal C^*F\mathcal C=G\) guarantees
its conservation in every state. In the canonical case,
\begin{equation}
 \mathcal T_8(G)=\frac{72G+R^*GR}{73}.
 \label{lib04:lector-72}
\end{equation}
This formula relates conservation and symmetry through a concrete
commutator. The observable that transports an operator
\(A\) of the domain exactly to the complete image is \(W_aAW_a^*\); its
ambient identity is \(W_aW_a^*\). The diagonal observable in
\eqref{lib04:lector-general} is a different choice, whose difference is
calculated by the same theorem.

\subsection{Reversible updating with incoming memory}
\label{lib04:reentrada}

Now let \(\mathcal B=I\), \(\mathcal C=U\), with \(U\)
unitary on the same space, and write
\begin{equation}
 A=\frac{\sqrt{a+1}}{\sqrt N}Q_-,\qquad
 D=\frac{\sqrt a}{\sqrt N}Q_+.
 \label{lib04:ad}
\end{equation}
The operators \(A,D\) are normal and commute with each other and with
their adjoints, since they are polynomials in \(U\) and \(U^*\).

\begin{theorem}[Unitary extension]
\label{lib04:teo-unitaria}
The update
\begin{equation}
 \mathscr R_a=
 \begin{pmatrix}A&-D^*\\D&A^*\end{pmatrix}:H\oplus H\to H\oplus H
 \label{lib04:unitaria}
\end{equation}
is unitary. Its first column is \(W_a\), and an input
\((v,m)\) preserves \(\|v\|^2+\|m\|^2\).
\end{theorem}

\begin{proof}
The balance gives \(A^*A+D^*D=I\). Normality yields
\(AA^*+DD^*=I\). The cross products of
\(\mathscr R_a^*\mathscr R_a\) and
\(\mathscr R_a\mathscr R_a^*\) vanish by the preceding
commutation relations; both diagonal blocks are the identity. The first column
reproduces \eqref{lib04:w}, and conservation of the joint norm
is the unitary identity applied to \((v,m)\).
\end{proof}

For data \(z=(z_1,z_2)\), the explicit inverse is
\begin{equation}
 v_{\mathrm{rec}}=A^*z_1+D^*z_2,\qquad
 m_{\mathrm{rec}}=-Dz_1+Az_2.
 \label{lib04:inversa-unitaria}
\end{equation}
The hypothesis of zero incoming memory is checked through
\(m_{\mathrm{rec}}=0\). If \(e\) is a perturbation of the data,
unitarity gives
\begin{equation}
 \|e\|^2=\|e_{\mathrm{rec}}\|^2+\|e_{\mathrm{mem}}\|^2.
 \label{lib04:descomposicion-ruido}
\end{equation}
The second component detects the part incompatible with that input.
An error belonging to \(\operatorname{im}W_a\) has zero memory
residual and changes the recovered input. Correction of the integer
record also uses its discrete separation distances.

\subsection{Causal reduction of the bilateral update}
\label{lib04:reduccion-causal}

The preceding unitary update realizes the transport of the state and its
memory in a common space. Observation of the first component admits an
equivalent description that explicitly retains the contribution of the
record. This description follows from the transport already constructed,
with the same coefficients and orientation: memory removed from the list
of variables reappears as dependence on earlier states.

Fix the operators \(A,D\) of \eqref{lib04:ad} and a realization
\(H\oplus H\). For \(n\geq0\), let
\begin{equation}
 v_{n+1}=Av_n-D^*m_n,\qquad
 m_{n+1}=Dv_n+A^*m_n.
 \label{lib04:recurrencia-memoria}
\end{equation}
The joint update is \(\mathscr R_a\) and preserves
\(\|v_n\|^2+\|m_n\|^2\).

\begin{proposition}[Exact elimination of the auxiliary record]
\label{lib04:prop-memoria-reducida}
For every \(n\geq0\), interpreting the empty sum as zero,
\begin{align}
 m_n&=(A^*)^n m_0+
       \sum_{k=0}^{n-1}(A^*)^{n-1-k}Dv_k,
       \label{lib04:memoria-explicita}\\
 v_{n+1}&=Av_n-D^*(A^*)^n m_0
       +\sum_{k=0}^{n-1}\mathscr K_{n-1-k}v_k,
       \label{lib04:evolucion-reducida}\\
 \mathscr K_j&=-D^*(A^*)^jD\quad(j\geq0).
       \label{lib04:nucleo-retardo}
\end{align}
The reduced equation, together with \eqref{lib04:memoria-explicita}, is
equivalent to the joint recurrence for the same data \((v_0,m_0)\).
\end{proposition}
\begin{proof}
The first identity holds for \(n=0\). If it holds for \(n\), the second
equation in \eqref{lib04:recurrencia-memoria} gives
\[
 m_{n+1}=(A^*)^{n+1}m_0+
 \sum_{k=0}^{n-1}(A^*)^{n-k}Dv_k+Dv_n,
\]
which is the identity at \(n+1\). Substitution at depth \(n\) into the
first equation proves \eqref{lib04:evolucion-reducida}. Conversely, if
\(v_n\) satisfies the reduced equation and \(m_n\) is defined by
\eqref{lib04:memoria-explicita}, the same difference of sums proves
\(m_{n+1}=Dv_n+A^*m_n\); substitution recovers the other equation.
\end{proof}

The kernel \(\mathscr K_j\) describes a transfer into the record,
\(j\) updates within it, and re-entry into the observed component.
The term \(-D^*(A^*)^nm_0\) preserves the initial memory. Omitting it
corresponds to the particular case \(m_0=0\), rather than an identity of
the general transport. Equation \eqref{lib04:evolucion-reducida} involves
only the current and earlier states. The causal character of this reduction
follows from the order of the joint recurrence.

The normalization of the balance also provides an explicit estimate.
Since \(Q_-=aI-U\), we have
\(\|A\|^2\leq r_a=(a+1)^3/((a+1)^3+a^3)<1\), whereas
\(D^*D\leq I\). Consequently,
\begin{equation}
 \|\mathscr K_j\|\leq r_a^{j/2},\qquad
 \sum_{j\geq0}\|\mathscr K_j\|
 \leq\frac{1}{1-\sqrt{r_a}}.
 \label{lib04:cota-nucleo-retardo}
\end{equation}
This bound controls the weight of memory in the reduced equation; the norm
of the joint state remains conserved by the unitary update.

The proposition is the temporal realization of the Schur reduction in
\cref{mem:prop-schur}. For the formal series
\(V(z)=\sum_{n\geq0}v_nz^n\) and
\(M(z)=\sum_{n\geq0}m_nz^n\), the recurrences yield
\(V=v_0+zAV-zD^*M\) and \(M=m_0+zDV+zA^*M\). Eliminating \(M\),
\begin{equation}
 V(z)=\bigl[I-zA+z^2D^*(I-zA^*)^{-1}D\bigr]^{-1}
       \bigl[v_0-zD^*(I-zA^*)^{-1}m_0\bigr].
 \label{lib04:schur-bilateral}
\end{equation}
The formal inverses exist because their constant terms are identities;
the state series and the joint resolvent also converge for \(|z|<1\).
The initial source and the effective operator transform together. A readout
that also acts on \(m_n\) retains its contribution through
\eqref{lib04:memoria-explicita}, in accordance with \eqref{mem:schur-lector}.

Re-entry here updates an incoming record. By contrast, the archiving policy
below retains each complement in a distinct coordinate. Each operation has
its own conservation law. The kernel \(\mathscr K_j\), Green propagators,
and a dimensional force have different types: the first eliminates a state
component, Green propagators solve source equations with support conditions, and
the force adds a dynamical realization and its units.


\subsection{Nonstationary archiving and uniform exhaustion of the terminal}
\label{lib04:archivo}

At each stage let \(H_n,H_{n+1}\) be Hilbert spaces and let
\(\mathcal B_n,\mathcal C_n:H_n\to H_{n+1}\) be isometries. The
parameter \(a>0\) remains fixed. Let \(A_n,D_n\) be the normalized
blocks of \eqref{lib04:w}; continuation proceeds through the first and
the second is recorded:
\begin{equation}
 R_0=I,\qquad R_{n+1}=A_nR_n,\qquad m_n=D_nR_nv.
 \label{lib04:politica-archivo}
\end{equation}
This policy has a different domain from memory reintroduction in
\eqref{lib04:unitaria}; each retains its own update.

\begin{lemma}[Uniform transfer]
\label{lib04:lem-tasa}
At every stage,
\begin{equation}
 \|A_n\|^2\leq r_a:={\frac{(a+1)^3}{(a+1)^3+a^3}}<1,
 \quad D_n^*D_n\geq(1-r_a)I,
 \quad\|R_n\|^2\leq r_a^n.
 \label{lib04:tasa}
\end{equation}
\end{lemma}

\begin{proof}
The triangle inequality gives
\(\|a\mathcal B_n-\mathcal C_n\|\leq a+1\). Normalization
yields the first bound. The difference
\(N-(a+1)^3=a^3>0\) proves \(r_a<1\). The local balance implies
\(D_n^*D_n=I-A_n^*A_n\), and multiplying the bounds for the
continued channels gives the estimate for \(R_n\).
\end{proof}

\begin{theorem}[Complete memory at arbitrary depth]
\label{lib04:teo-archivo}
The finite map
\begin{equation}
 Z_Nv=(R_Nv,D_0v,D_1R_1v,\ldots,D_{N-1}R_{N-1}v)
 \label{lib04:zn}
\end{equation}
is isometric. The infinite archive
\begin{equation}
 Z_\infty v=(D_nR_nv)_{n\geq0}
 \label{lib04:zinfty}
\end{equation}
belongs to the Hilbert direct sum of the output spaces and satisfies
\begin{equation}
 \|v\|^2=\sum_{n\geq0}\|m_n\|^2,\qquad
 v=\sum_{n\geq0}R_n^*D_n^*m_n,
 \qquad Z_\infty^*Z_\infty=I.
 \label{lib04:recuperacion-infinita}
\end{equation}
The inverse series converges in norm and the adjoint has norm one.
\end{theorem}

\begin{proof}
The identity \(A_n^*A_n+D_n^*D_n=I\), conjugated by \(R_n\),
gives
\[
 R_n^*R_n-R_{n+1}^*R_{n+1}=R_n^*D_n^*D_nR_n.
\]
Summing from zero to \(N-1\) yields
\begin{equation}
 I=R_N^*R_N+\sum_{n<N}R_n^*D_n^*D_nR_n.
 \label{lib04:telescopia}
\end{equation}
This is the Gram identity for \(Z_N\). By
\eqref{lib04:tasa}, the first term tends to zero in operator
norm. The limit identity proves square summability of the
archive and its isometric property. The adjoint of this isometry is bounded;
the truncations of any square-summable archive converge in the Hilbert
direct sum, and their images under the adjoint prove convergence of the
series in \eqref{lib04:recuperacion-infinita}. Applying it to the exact
data recovers \(v\).
\end{proof}

With only the first \(N\) complements, the truncated adjoint has
the exact residual
\begin{equation}
 v-\sum_{n<N}R_n^*D_n^*m_n=R_N^*R_Nv,
 \qquad
 \left\|v-\sum_{n<N}R_n^*D_n^*m_n\right\|
 \leq r_a^N\|v\|.
 \label{lib04:sesgo-truncado}
\end{equation}
If the joint data error is at most \(\varepsilon\), the
total bound is \(\varepsilon+r_a^N\|v\|\). With \(Z_N\), including
its terminal, or with \(Z_\infty\), the adjoint recovers the input with error
at most \(\varepsilon\).

The preceding policy exhausts the terminal through the proven uniform bound.
The earlier evolutions with a positive terminal
\(A_\infty\neq0\) retain the component
\(A_\infty^{1/2}v\) in their archive. Arbitrary depth and elimination of the
terminal are therefore different conclusions, each retaining its
respective hypotheses.

\begin{proposition}[Inverse of finite memories without the terminal]
\label{lib04:prop-inversa-finita}
For \(N\geq1\), let \(M_Nv=(D_nR_nv)_{n<N}\). Then
\begin{equation}
 M_N^*M_N=I-R_N^*R_N\geq(1-r_a^N)I,
 \label{lib04:gram-finito}
\end{equation}
and the least-squares inverse
\begin{equation}
 L_N=(I-R_N^*R_N)^{-1}M_N^*,\qquad
 L_NM_N=I,\qquad
 \|L_N\|\leq(1-r_a^N)^{-1/2}
 \label{lib04:inversa-finita}
\end{equation}
recovers the input exactly from exact data.
\end{proposition}

\begin{proof}
The Gram identity is \eqref{lib04:telescopia}; the bound follows from
\eqref{lib04:tasa}. The positive operator is invertible by its strict
lower bound. Composition proves \(L_NM_N=I\), and
\(L_NL_N^*=(M_N^*M_N)^{-1}\) proves the stated norm bound.
\end{proof}

This finite inverse removes the bias in
\eqref{lib04:sesgo-truncado}; its noise amplification is
specified by \eqref{lib04:inversa-finita}. The truncated adjoint,
finite inverse, and adjoint of the complete archive are thus three
operators with distinct behavior.

For \(a=8\), the general bound is \(r_a=729/1241\). In the case
\(\mathcal B=I\), \(\mathcal C=S^4\), with \(S\) the cycle of
twelve positions, the fixed-space projector
\(P_0=(I+S^4+S^8)/3\) gives
\begin{equation}
 A^*A=\frac{a+1}{(2a+1)p}
       \bigl((a-1)^2P_0+p(I-P_0)\bigr).
 \label{lib04:gram-tasa-ternaria}
\end{equation}
Since \(p-(a-1)^2=3a>0\), its norm is
\((a+1)/(2a+1)\). At \(a=8\), this is \(9/17\); the eigenvalue
of the fixed sector is \(441/1241\). These rates belong to the
declared transports and policy.

If both channels are refined in a tree, every complete level is also
isometric by composition. Summing the norms of all complete layers
would count the same initial norm repeatedly. Conservation
is formulated through the refinement maps between levels or
through the archive of complements in \eqref{lib04:zinfty}.

\subsection{Geometric forms and conservation under transport}
\label{lib04:formas}

Let \(b,b'\) be nondegenerate symmetric or alternating forms on
the corresponding spaces, and suppose that
\(\mathcal B^*b'\mathcal B=\mathcal C^*b'\mathcal C=b\).
The same bilinear expansion proves
\begin{equation}
 (a+1)Q_-^*b'Q_-+aQ_+^*b'Q_+=Nb.
 \label{lib04:balance-formas}
\end{equation}
Indeed, the mixed terms
\(\mathcal B^*b'\mathcal C+\mathcal C^*b'\mathcal B\) cancel
with the same weights; the pure terms sum to \(Nb\).
The transported form is thus preserved without requiring positivity.
The contraction and noise bounds of the preceding subsection, by
contrast, use positive Hilbert norms. A Lorentzian or alternating
form retains \eqref{lib04:balance-formas}, with its own geometric
interpretations, and does not replace positivity in those
inequalities.

\subsection{Channel correlation and an inverse test of a reading}
\label{lib04:correlacion}

For \(v\neq0\), define
\begin{equation}
 \chi(v)=\frac{\operatorname{Re}\langle\mathcal Bv,\mathcal Cv\rangle}
                      {\|v\|^2},\qquad -1\leq\chi(v)\leq1.
 \label{lib04:chi}
\end{equation}
The bound follows from Cauchy--Schwarz applied to the two isometric images.
The normalized readings are
\begin{equation}
 s_-=a^2+1-2a\chi,\qquad
 s_+=(a+1)^2+1+2(a+1)\chi.
 \label{lib04:s-pm}
\end{equation}
For \(a=8\),
\begin{equation}
 s_-=65-16\chi,\qquad s_+=82+18\chi,\qquad
 9s_-+8s_+=1241.
 \label{lib04:s-canonico}
\end{equation}
The two individual norms contain a common real correlation. The
balance is an exact relation between them, not two independent
arithmetic determinations of a constant.

Using the previous output \(\alpha_A\) in the inverse test
\(s_-=10^4\alpha_A\) determines the complementary reading
\begin{equation}
 s_+=73+\frac98(73-10^4\alpha_A)
       \simeq73{,}029783595557.
 \label{lib04:ensayo-alpha}
\end{equation}
The equality follows by substitution into
\eqref{lib04:s-canonico}. The weighted average
\((9s_-+8s_+)/17\) remains exactly \(73\).
The comparison preserves the prior genealogy of \(\alpha_A\): it is
a condition of the inverse test, not another independent producer of
that output. A unitary change can alter \(\chi\) and both
readings while preserving the balance, even if a particular ternary-order
relation changes.

\subsection{Cyclic frame and transport of responses}
\label{lib04:marco}

On the carrier of the generated record, let
\(H=\mathbb R^{12}\), \((Sx)_i=x_{i+1}\), with cyclic indices,
and
\begin{equation}
 O_K=(K,SK,\ldots,S^{11}K),\qquad
 589609I\leq O_K^*O_K\leq39225169I.
 \label{lib04:gram-k}
\end{equation}
These are the bounds for the cyclic frame of the canonical record,
proved in \cref{lib01:marco-ciclico}.
Let \(Z\) be one of the complete isometries
\(Z_N,Z_\infty\), or a composition with the preceding archives
that retains their terminals and complements.

\begin{theorem}[Preservation of frame conditioning]
\label{lib04:teo-marco}
If \(x=O_K\beta\), the datum \(z=Zx\) determines
\begin{equation}
 \beta=O_K^{-1}Z^*z,\qquad
 \|\delta\beta\|\leq\frac{\|\delta z\|}{\sqrt{589609}}.
 \label{lib04:inversa-marco}
\end{equation}
The twelve columns of \(ZO_K\) form a basis of its image, even
when the archive has infinitely many coordinates.
\end{theorem}

\begin{proof}
The equality \(Z^*Z=I\) gives
\((ZO_K)^*(ZO_K)=O_K^*O_K\). The Gram bounds are preserved, and
\(O_K\) is invertible by the strict lower bound. Composing their
inverses proves recovery, and the norm bound
\(\|O_K^{-1}\|\leq589609^{-1/2}\), together with
\(\|Z^*\|=1\), gives stability. The number of columns and their
independence determine the dimension of the image, without identifying it
with the entire storage space.
\end{proof}

For an operator \(H_0\) commuting with \(S\), the complete vector
response \(y=H_0K\) determines
\begin{equation}
 H_0=O_yO_K^{-1}.
 \label{lib04:identificacion-filtro}
\end{equation}
Indeed, commutation implies
\(H_0O_K=(y,Sy,\ldots,S^{11}y)=O_y\). If only
\(z=Zy+e\) is received, define \(\widehat y=Z^*z\) and
\(\widehat H_0=O_{\widehat y}O_K^{-1}\). Then
\begin{equation}
 \|\widehat H_0-H_0\|
 \leq\sqrt{\frac{12}{589609}}\,\|e\|.
 \label{lib04:error-filtro}
\end{equation}
The orbital matrix of \(\delta y\) has twelve columns of norm
\(\|\delta y\|\), so its operator norm is at most
\(\sqrt{12}\|\delta y\|\). This and
\eqref{lib04:gram-k} prove \eqref{lib04:error-filtro}.
For a general operator, the protocol requires the twelve responses
\(H_0S^jK\); a response in the cyclic case means its twelve
components, not a scalar measurement.

In the transported image, the shift is
\(S_Z=ZSZ^*\), and satisfies \(S_Z^jZ=ZS^j\). This equality
determines phase advance in the archive; a shift of slots in the
ambient space represents that advance only when it coincides with
the transported operator.

\subsection{Positional reading and scalar precision}
\label{lib04:kappa}

Fix \(B_0=1000\), \(d=B_0^{12}-1\), and the row
\begin{equation}
 \ell=\frac{(B_0^{11},B_0^{10},\ldots,1)}{d}.
 \label{lib04:fila-posicional}
\end{equation}
Its linear extension acts on \(H\), and on the domain of words
\(\kappa(K)=\ell K\). The representative of the transported reader is
\(Z\ell^*\), so
\begin{equation}
 \kappa(K)=\langle Z\ell^*,ZK\rangle,\qquad
 |\delta\kappa|\leq\|\ell\|\,\|e\|.
 \label{lib04:lectura-kappa}
\end{equation}
The geometric sum gives
\begin{equation}
 \|\ell\|^2
 =\frac{B_0^{12}+1}{(B_0^2-1)(B_0^{12}-1)}.
 \label{lib04:norma-fila}
\end{equation}
Indeed, the numerator of the sum of squares is
\((B_0^{24}-1)/(B_0^2-1)\), which, when divided by
\((B_0^{12}-1)^2\), yields \eqref{lib04:norma-fila}.

With only \(N\) complements and the truncated adjoint, there is an
additional error of at most \(\|\ell\|r_a^N\|K\|\). The finite inverse
\eqref{lib04:inversa-finita} removes that bias and retains its explicit
noise-amplification bound.

On the domain \(\{0,\ldots,999\}^{12}\), the exact reading
recovers \(N_K=d\kappa\) and its twelve blocks through successive
Euclidean divisions. If
\begin{equation}
 |\widetilde\kappa-\kappa|<\frac1{2d},
 \label{lib04:radio-escalar}
\end{equation}
rounding \(d\widetilde\kappa\) to the nearest integer
recovers \(N_K\): its distance from the true integer is less than
half a unit. The base, length, origin, and orientation accompany
that representation. The scalar threshold in
\eqref{lib04:radio-escalar} is distinct from the vector bounds
for memory records.

\subsection{Incidence, charge, and exact recovery of the record}
\label{lib04:incidencia}

Let \(\mathbf1=(1,\ldots,1)^T\) and
\(\Pi=I-\mathbf1\mathbf1^*/12\). The incidence isometry of
\cref{nuc05:isometria} is
\begin{equation}
 Fv=(\mu(v),\mathcal I\Pi v/6),\qquad
 \mu(v)=\frac{\mathbf1^*v}{12},\qquad
 \|(\mu,y)\|_Y^2=12\mu^2+\|y\|^2.
 \label{lib04:f-incidencia}
\end{equation}
Its codomain is its declared image. The Gram identity
\(\mathcal I^*\mathcal I=36I+30\mathbf1\mathbf1^*\) gives
\(\|\mathcal I\Pi v\|^2=36\|\Pi v\|^2\), and
\(12\mu(v)^2+\|\Pi v\|^2=\|v\|^2\), hence \(F^*F=I\).
Thus, from \(z=Zv+e\),
\begin{equation}
 \widehat{Fv}=FZ^*z,\qquad \|FZ^*e\|_Y\leq\|e\|.
 \label{lib04:error-incidencia}
\end{equation}
The unnormalized incidence \(\mathcal I\Pi v\) has error at most
\(6\|e\|\); the signed record \(H_{12}v\), whose normalization
satisfies \(H_{12}^*H_{12}=4I\), has error at most
\(2\|e\|\).

The direction \(u=P_3\Pi K\) and the projector
\(P_{10}=\Pi-uu^*/\|u\|^2\), with \(u\neq0\), are recovered
with their previously defined readers and chart. The absolute selection
\(\{j:K_j\geq729\}\) has a discontinuity at \(K_4=729\).
Integer recovery is performed before applying that threshold.

\begin{proposition}[Decoding after the complete archive]
\label{lib04:prop-redondeo}
Let \(K\) be integral and let \(z=ZK+e\), with \(Z^*Z=I\). If
\(\|e\|<1/2\), rounding each coordinate of \(Z^*z\) recovers
\(K\) exactly. Composition then restores \(\kappa\), the signatures,
and the geometric readers defined on that record.
\end{proposition}

\begin{proof}
Contractivity of \(Z^*\) gives
\(\|Z^*z-K\|\leq\|e\|<1/2\). Each coordinate error is
bounded by the joint norm, so the unique integer coordinate at
distance less than half a unit is the original one. Once
\(K\) has been recovered, evaluating any of its exact readers
restores its value.
\end{proof}

In particular, the support \(\{4,6,9,10\}\) is recovered, including
its boundary component. The other marks involved in an exceptional
flag are retained by their own channels: recovery of
\(K\) corrects exactly the information for which it supplies
sufficient arguments.

\subsection{Preservation of the conditioning of two memories}
\label{lib04:dos-memorias}

Retain the preceding readers
\begin{equation}
 T_j=\frac{8I+S^j}{9},\qquad
 D_j=\frac{\sqrt8}{9}(S^j-I),\qquad
 Mv=(D_3v,D_4T_3v).
 \label{lib04:m-antecedente}
\end{equation}
The \(D_j\) in this equation are the preceding complements;
the second bilateral block \(D\) of \eqref{lib04:ad} retains
its own definition. One has \(\ker M=\mathbb R\mathbf1\)
and an inverse \(L_M\) of norm \(9/4\) on the centered sector,
as proved in \cref{nuclear:11}.

Let \(E=W_a\oplus W_a\) on the space of the two memories.
Its isometric property implies
\begin{equation}
 (EM)^*(EM)=M^*M,\qquad
 \ker(EM)=\mathbb R\mathbf1.
 \label{lib04:gram-m-preservado}
\end{equation}
For \(z=EMv+e\) and exact charge
\(q=\mathbf1^*v\), the reconstruction is
\begin{equation}
 \widehat v=\frac q{12}\mathbf1+L_ME^*z,\qquad
 \|\widehat v-v\|\leq\frac94\|e\|.
 \label{lib04:inversa-m}
\end{equation}
The proof consists in applying \(E^*E=I\), using
\(L_MM=\Pi\) and \(\|E^*\|=1\). If the charge has error
\(\delta q\), the uniform and centered sectors are orthogonal, and
\begin{equation}
 \|\widehat v-v\|^2
 \leq\frac{|\delta q|^2}{12}+\frac{81}{16}\|e\|^2.
 \label{lib04:error-carga}
\end{equation}

The distances between all words in the integer image are preserved
exactly by \(E\). Therefore, the recovery radii of
\cref{lib02:radios-memoria} remain
\begin{equation}
 r_0=\frac{2\sqrt{146}}{81}\quad\hbox{without charge},\qquad
 r_q=\frac{4\sqrt{73}}{81}\quad\hbox{with exact charge}.
 \label{lib04:radios}
\end{equation}
Within the first radius, \(\Pi K\) is recovered, and with it the
direction and its projector whenever they are defined. Within the second,
with \(q\) retained, one recovers \(K\) and the absolute threshold
selector. Subsequent composition with another complete isometry, finite
or infinite, preserves the same guarantees: its adjoint is applied before
the decoder. The radius is always expressed in the joint data
norm with the fixed normalization.

This composition redistributes the information in \(M\) without degrading
its conditioning or creating the uniform channel that \(M\) eliminates.
Gram preservation distinguishes this property from a numerical
improvement obtained by changing the reader or the noise scale.

\subsection{Orbital norm and polar normalization of the frame}
\label{lib04:normalizacion}

The preceding inversion satisfies
\(K=H_{12}U_{\mathrm{sgn}}/4\),
\(H_{12}^*H_{12}=4I\). Therefore
\begin{equation}
 \|K\|^2=\frac14\|U_{\mathrm{sgn}}\|^2=4251257,
 \qquad \|U_{\mathrm{sgn}}\|^2=17005028.
 \label{lib04:norma-k}
\end{equation}
Each shift \(S^j\) is orthogonal and preserves
\(\nu=\sqrt{4251257}\). Hence
\begin{equation}
 \operatorname{tr}(O_K^*O_K)=12\nu^2=51015084.
 \label{lib04:traza-marco}
\end{equation}
Division by \(\nu\) normalizes each orbital column to norm
one and preserves its order. It is uniform over the orbit and under its
isometric transports. Different records retain their own
norms, and the normalized columns still have a nonscalar
Gram matrix: its eigenvalues are the preceding ones divided by
\(4251257\).

Operator normalization uses the complete Gram matrix:
\begin{equation}
 G_K=O_K^*O_K,\qquad
 F_K=O_KG_K^{-1/2},\qquad
 F_K^*F_K=I.
 \label{lib04:polar}
\end{equation}
The inverse square root exists by \eqref{lib04:gram-k}. The equality is
verified by multiplying the Gram matrix on both sides by \(G_K^{-1/2}\).
Thus \(ZF_K\) is isometric at any depth.
The complete recovery of the frame is
\begin{equation}
 O_K=F_KG_K^{1/2}.
 \label{lib04:polar-inversa}
\end{equation}
The pair \((F_K,G_K)\) preserves the complete information.
The unitary factor retains an oriented orthonormal frame in the fixed
charts, whereas the amplitudes also reside in
\(G_K\). Multiplying \(K\) by a positive scalar preserves
\(F_K\) and changes \(G_K\); \(K\) and \(-K\) share a
Gram matrix and have opposite polar factors, with different linear readings
\(\kappa\). Each example shows why an isolated factor cannot
replace the pair.

Shifting the origin preserves \(\nu\), but in the fixed positional
chart
\begin{equation}
 \kappa(SK)=1000\kappa(K)-K_1.
 \label{lib04:cambio-origen}
\end{equation}
This follows by shifting the positional numerator and using
\(d=1000^{12}-1\). To preserve the same reading under a change
of chart, the row \(\ell\) is also transported by the inverse
of that change. A general unitary chart can leave the integer-block
lattice; decoding by Euclidean division is performed
in the recovered block chart.

\subsection{An exact example of transduction and error correction}
\label{lib04:ejemplo}

Take \(a=8\), \(U=S^4\), and the canonical record
\begin{equation}
 \begin{split}
 K=(&234,543,140,729,659,824,\\
    &621,58,914,794,146,601).
 \end{split}
 \label{lib04:k-ejemplo}
\end{equation}
Denoting its two blocks of six coordinates by \((v)_{1:6}\) and
\((v)_{7:12}\), the unnormalized readings are
\begin{align}
 (Q_-K)_{1:6}&=(1213,3520,499,5774,4358,5798),\nonumber\\
 (Q_-K)_{7:12}&=(4822,-137,7078,5809,1028,4079),\nonumber\\
 (Q_+K)_{1:6}&=(2765,5711,1881,6619,6845,8210),\nonumber\\
 (Q_+K)_{7:12}&=(5735,1123,8460,7689,1454,6138).
 \label{lib04:canales-ejemplo}
\end{align}
Their sum, divided by \(17\), recovers every block of \(K\).
The matrix identity
\(9Q_-^*Q_-+8Q_+^*Q_+=1241I\) precedes its evaluation on the
particular vector.

To introduce a controlled error, the two memories are rationalized
as \(z_0=MK/\sqrt8\), and \(1/10\) is added to the first
coordinate. Each block of twelve components is then transported
through the two bilateral channels. The sum inverse recovers the
perturbed memory data, and the charge decoder receives
\(q=6263\) together with those data, without receiving the record sought.
The norm of the error in the original memories is
\(\sqrt8/10<r_q\). Isometric normalization preserves that norm;
the unnormalized channels merely facilitate exact rational calculation.

In this instance, there are \(924\) candidates from floor-and-ceiling
reconstruction that satisfy the charge. The minimum rationalized squared
distance is \(1/100\), and the decoder recovers
\eqref{lib04:k-ejemplo}. Uniqueness follows from the strict radius
\eqref{lib04:radios}, not merely from the size of the enumeration.
The recovered arithmetic reading is
\begin{equation}
 \kappa=
 \frac{234543140729659824621058914794146601}
      {999999999999999999999999999999999999},
 \label{lib04:kappa-ejemplo}
\end{equation}
and the incidence selection is again \(\{4,6,9,10\}\), with
the same threshold \(729\). The example introduces noise before
transport. The adjoint contractivity proofs also cover arbitrary
subsequent errors whose joint norm lies within the
declared radius.

The combined operational capability is to transport and recover a generated
record, preserve the conditioning of its responses, and restore
its arithmetic and incidence readings before applying discontinuous
selections. Its effective premises are the prior production of
\(K\), the invertible cyclic frame, the bilateral balance, the
discrete separation of memories, and the archiving theorem. The
twelve-component record retains the scope of its readers;
the complete enriched history retains its other records and
dependencies.
